\subsection{Objectifs et contraintes initiales}


La dernière étape de notre projet consistait à concevoir une enveloppe physique pour protéger nos cartes électroniques. L'objectif était de créer un boîtier sur mesure, imprimable en 3D, capable d'accueillir notre circuit imprimé tout en laissant un accès aux connectiques.
Cependant, cette phase comportait une contrainte : nous étions dans l'obligation d'attendre la fin de la conception du PCB sous KiCad. En effet, la moindre modification du placement des composants ou des trous de fixation sur la carte avait un impact direct sur la géométrie de notre futur boîtier. Il en allait de même pour la taille et la hauteur des composants, pour lesquelles nous devions consulter les datasheets et effectuer des mesures manuelles précises afin d'obtenir l'espacement le plus optimisé possible.


\subsection{Découverte de l'outil et problématiques de prise en main}


Pour réaliser cette modélisation, nous avons choisi d'utiliser le logiciel Autodesk Fusion 360. La modélisation 3D était une étape relativement nouvelle pour nous, car personne n'avait d'expérience préalable dans ce domaine. 
La prise en main du logiciel s'est avérée complexe. Face au peu de tutoriels disponibles pour nous accompagner spécifiquement sur l'intégration d'une carte électronique dans un boîtier imprimable, nous avons procédé par itération. Nous avons donc décidé de réaliser une "boîte test" pour commencer. Ce premier essai nous a permis de comprendre le fonctionnement de l'espace de travail, la gestion de l'historique de conception, et l'utilisation des outils de base tels que les esquisses 2D et les extrusions.


\subsection{Schéma général de conception : La méthode "Top-Down"}
Une fois familiarisés avec Fusion 360, nous avons pu mettre en place une méthodologie plus rigoureuse appelée conception "Top-Down". Au lieu de modéliser le boîtier à l'aveugle, nous avons importé le modèle 3D de notre carte électronique (le fichier .STEP généré par KiCad) directement au centre de notre espace de travail. 


Voici les fonctions clés que nous avons utilisées pour construire le boîtier autour d'elle :
L'outil Projeter (Raccourci P) : Cet outil a été fondamental dans notre conception. Il nous a permis de "calquer" le contour exact de la carte électronique sur le fond ou le couvercle de notre boîtier. Cela nous a permis de garantir que les trous de fixation pour les vis soient positionnés au millimètre près pour l’assemblage et l'impression 3D.



L'outil Décalage (Offset) : À partir du contour projeté, nous avons défini l'épaisseur des murs extérieurs. Nous avons opté pour une épaisseur de 2.5 mm. Cette valeur a été choisie en fonction des recherches que nous avons effectuées sur le rechargement sans fil mais aussi sur les caractéristiques techniques des imprimantes 3D qui étaient à notre disposition.
\subsection{Assemblage et système de fermeture}


La problématique majeure de notre modélisation était de concevoir un système permettant de maintenir la carte fermement au fond du boîtier, tout en ayant un couvercle vissable et démontable.
Pour ce faire, nous avons modélisé des pylônes de fixation directement fusionnés aux quatre coins intérieurs de la boîte. L'outil d’Extrusion a été paramétré sur l'opération "Joindre" pour que ces piliers de 8 mm de diamètre fassent corps avec les murs.
Ensuite, nous devions percer ces piliers pour le passage des vis. Nous avons utilisé de nouveau l'outil Extrusion, mais cette fois-ci avec l'opération "Couper", pour vider le centre des pylônes avec un diamètre de 3.2 mm. Ce diamètre correspond aux vis M3 standard que nous avons utilisées.
À la base, nous voulions utiliser des vis en plastique pour la fixation, afin d'éviter tout problème de court-circuit sur les cartes électroniques. Cependant, avec l'équipement à notre disposition, nous ne pouvions pas réaliser les rainures adaptées pour ce type de vis. Nous nous sommes donc rabattus sur des vis métalliques. Il reste néanmoins possible d'ajouter un plot anti-vibration pour sécuriser davantage l'installation.



Nous avons aussi dû prendre en compte les conditions d'utilisation finales : le boîtier manipulé par l'ouvrier ne possède pas de système de fermeture à vis avec une membrane pour l'étanchéité, alors que les couvercles de nos autres boîtiers sont fermés avec des vis pour en faciliter la maintenance. 
Enfin, avant de lancer la fabrication, nous avons utilisé le logiciel de découpe de Bambu Lab. Cela nous a permis de visualiser finement nos boîtes couche par couche et de vérifier tous les paramètres avant l'impression 3D. 
